%----------------------------------
\section{Peer Review}
\label{sec:Peer Review}
%----------------------------------

\noindent\emph{Peer reviewed: Abhishek Mathur --- ``Seeing Around the Price Spike: Predictive
Modelling for Customer \& Price Risk at Alinta Energy'' (Individual Task 1: Part 2 video).}

\subsection*{Soundness and Significance}

The presentation frames a genuinely significant commercial problem well: wholesale price spikes
and heterogeneous household usage as two distinct risk types for an energy retailer. The
trade-off framing on the results slide --- one model catching 93\% of spikes but with more false
alarms, versus 82\% precision when flagged --- is a sound way to communicate a precision/recall
trade-off without naming it as such, and correctly avoids declaring a false ``winner.'' Flagging
that the household model favoured simplicity because only 25 households were available is an
honest, important caveat about sample size that a less careful presenter might have glossed over.

\subsection*{Executive-Level Insight}

This is the presentation's strongest area. The ``over-sensitive smoke detector'' analogy for
false alarms, and the closing ``market data shows WHEN, household data shows WHO'' framing, are
exactly the kind of memorable, jargon-free translations an executive audience needs. The two
datasets are positioned as complementary rather than competing, which gives the pitch a clear
takeaway rather than leaving it as two disconnected technical results.

\subsection*{Presentation Quality}

Visual design is clean and consistent --- icon-led cards, one message per slide, consistent
colour coding. Two issues affect delivery, though. First, pacing is uneven: the title slide alone
runs close to 30 seconds before any content appears, while the results slide consumes roughly a
third of the total three-minute runtime, leaving the closing synthesis slide comparatively
rushed. Second, the webcam overlay in the bottom-right corner covers part of the on-screen text
on at least two slides (the household dataset size, and the second model's comparison figure), so
a viewer relying on the visuals alone loses some of the numbers being discussed.

\subsection*{Suggestions for Improvement}

\begin{enumerate}
\item Split the results slide into two (spike detection, then household behaviour) so each gets
proper visual space and pacing evens out.
\item Reposition or shrink the webcam bubble so it never sits over data.
\item Consider naming the two modelling approaches in general terms (e.g.\ ``a flexible model''
vs.\ ``a simpler baseline'') so claims are traceable without requiring jargon.
\item The ending states the insight well but could close with one explicit recommended action for
the business, to fully land the executive pitch.
\end{enumerate}

\noindent \textbf{Overall:} a well-structured, appropriately pitched narrative let down slightly
by pacing and a technical overlay issue.

\wordcount{398}
